% ECONOMICS_PROBLEM: {"id": "L22-89", "title": "Vertical integration trade-offs", "jel_code": "L22", "source_id": "OP-0396"}
% !TeX program = xelatex
\documentclass[11pt,a4paper]{article}
\usepackage[margin=23mm]{geometry}
\usepackage{fontspec}
\setmainfont{Latin Modern Roman}
\usepackage{amsmath,amssymb,mathtools}
\usepackage{xcolor,enumitem,booktabs,longtable,array,xurl,hyperref}
\definecolor{muted}{HTML}{53636C}
\hypersetup{colorlinks=true,urlcolor=blue,linkcolor=blue}
\setlength{\parindent}{0pt}
\setlength{\parskip}{5pt}
\newcommand{\R}{\mathbb R}
\newcommand{\E}{\mathbb E}
\newcommand{\N}{\mathbb N}
\newcommand{\ind}{\mathbf 1}
\newcommand{\Var}{\operatorname{Var}}
\newcommand{\Cov}{\operatorname{Cov}}
\newcommand{\supp}{\operatorname{supp}}
\DeclareMathOperator*{\argmin}{arg\,min}
\DeclareMathOperator*{\argmax}{arg\,max}
\newcommand{\entrymeta}[3]{{\sffamily\small\color{muted}\textbf{\texttt{#1}}\quad|\quad #2\par #3\par}}
\newcommand{\Problemref}[1]{\texttt{#1}}
\newcommand{\JELref}[2]{\texttt{#2} (JEL \texttt{#1})}
\begin{document}
\section*{Vertical integration trade-offs}
\textbf{Problem ID:} \texttt{L22-89}\quad \textbf{JEL:} \texttt{L22}\par
% BEGIN ECONOMICS STATEMENT
\label{problem:OP-0396}
{\sffamily\small\color{muted}Original subject group: Industrial organization and competition\par}
\label{definitions:problem:30007678}
\textbf{Audit classification:} Published research agenda. The cited chapter supports a broader modeling agenda. No result is universal over arbitrary unspecified bargaining protocols, and equilibrium selection or strong set comparisons must be explicit.
\subsection*{Definitions and precise research target}
Consider an oligopolistic supply chain with a finite upstream-firm set \(U\), finite downstream-firm set \(D\), and consumers with a specified demand system. Firms have positive differentiable production costs and may bargain over contracts combining per-unit prices and fixed transfers. Let \(\theta\) collect demand, costs, bargaining weights, entry costs and the information available during negotiation. In regime \(a=1\), one specified upstream-downstream pair integrates; in regime \(a=0\), all firms remain independent. In either regime firms may change contracting partners, invest and enter, subject to the same technological constraints. Let \(\mathcal E_a(\theta)\) be the set of equilibria under a stated bargaining protocol. Restrict parameter comparisons to cases where both equilibrium sets are nonempty, and define \(W(e)\) as consumer surplus plus firm profits, with profits measured net of real production, investment and entry costs in equilibrium \(e\). Transfers between firms cancel from this welfare measure. Characterize the parameter regions in which every \(e_1\in\mathcal E_1(\theta)\) and every \(e_0\in\mathcal E_0(\theta)\) satisfy \(W(e_1)\geq W(e_0)\). When this strong comparison fails, provide welfare ranges under an explicit equilibrium-selection rule. The task isolates efficiencies, foreclosure and partner rematching within a declared model class. It makes no universal assertion that integration is beneficial or harmful across all possible contracts and information structures.

\textbf{Definitions, assumptions and mathematical qualifications.} A complete input economy specifies finite firm sets, consumers' quasi-linear utility and outside option, feasible contracts, investment and entry action sets, stage order, beliefs, cost functions, shock law and a named equilibrium concept for the resulting game. For an equilibrium \(e\), \(W(e)\) is aggregate gross consumption utility less real costs; domestic payments cancel. Restrict to economies with nonempty equilibrium sets and bounded measurable welfare. Strong welfare dominance means \(\inf_{e_1\in\mathcal E_1(\theta)}W(e_1)\geq\sup_{e_0\in\mathcal E_0(\theta)}W(e_0)\). If using a selection rule, give a probability kernel \(\kappa_a(de\mid\theta)\) supported on \(\mathcal E_a(\theta)\) and compare \(\int W(e)\kappa_a(de\mid\theta)\).
\subsection*{Variants and assumption changes}
\begin{enumerate}
\item \textbf{Editorial, fixed links.} Hold contracting partners and investments fixed and permit only linear wholesale prices. Compare this simpler double-markup benchmark with the integrated regime under the stated demand model.
\item \textbf{Editorial, dynamic links.} Permit partner rematching, private cost information and persistent investments. Fix a finite horizon or discount factor, and identify welfare sign ranges using contract and investment observations.
\end{enumerate}
\subsection*{References and source audit}
\textbf{Checked source.} 14 September 2021 author draft, Section 5, printed pp. 58--59. The conclusion identifies richer contracts, endogenous partners, information and dynamics as research directions. Full primary text retrieved. \url{https://robinlee.sites.fas.harvard.edu/papers/HandbookIO_Vol4_VerticalMarkets.pdf}.
\begin{enumerate}\raggedright
\item\hypertarget{definitions:source:30007678:1}{} Robin S. Lee; Michael D. Whinston; Ali Yurukoglu. 2021. \emph{Structural Empirical Analysis of Contracting in Vertical Markets}. 14 September 2021 author draft, Section 5, printed pp. 58--59.
\url{https://robinlee.sites.fas.harvard.edu/papers/HandbookIO_Vol4_VerticalMarkets.pdf}.
DOI: \url{https://doi.org/10.3386/w29282}.
\end{enumerate}

\subsection*{Common conventions: Source mathematical conventions}

These conventions apply to each entry unless explicitly replaced.

\textbf{Models and identification.} A primitive model \(m\) specifies all probability laws, feasible choices, information, equilibrium and measurement rules used by its target. The admissible class \(\mathcal M\) is fixed independently of outcomes. If \(P_O(m)\) is its observable probability law and \(\tau(m)\) the target, its identified set at observable law \(P\) is
\[
 \mathcal I_\tau(P)=\{\tau(m):m\in\mathcal M,\ P_O(m)=P\}.
\]
Point identification means that this set is a singleton. An empty set signals incompatibility of data and assumptions. Coordinate bounds need not describe the joint set. Bounds are sharp only if every retained target is attainable or arbitrarily approachable by compatible models. Finite bounds require explicit restrictions on unobserved quantities; unrestricted confounding, hidden wealth, extrapolation or arbitrary equilibrium selection can yield uninformative sets.

\textbf{Causal interventions.} A potential outcome \(Y_i(p)\) is the outcome under a complete policy regime \(p\), including timing, financing, information and market spillovers. The target population and units are fixed before treatment. With interference, use the complete assignment vector or a justified exposure mapping; individual no-interference notation alone cannot identify an economy-wide intervention. A difference of mean potential outcomes does not require the cross-world joint distribution. Conditioning on post-treatment migration, survival, enrollment or employment changes the estimand and needs separate justification.

\textbf{Statistical statements.} An estimator, confidence set, forecast or test specifies a sampling law, model class and loss/coverage criterion. Uniform \((1-\alpha)\)-coverage means \(\inf_{m\in\mathcal M}\Pr_m\{\tau(m)\in C_n\}\geq1-\alpha\), or an explicitly asymptotic version. The whole parameter space trivially has coverage, so informativeness requires a stated width, risk or power objective. A forecasting improvement is relative to a named benchmark, information set and evaluation population.

\textbf{Equilibrium and optimization.} An equilibrium comprises feasible plans, optimality, beliefs where needed, budget constraints and all market/resource clearing. The primitives or a stated benchmark define each correspondence \(\mathcal E(p;m)\). Nonemptiness is assumed or proved, never inferred just from compact policy sets. A selected-equilibrium target uses a declared selection rule; a robust target quantifies over all permitted equilibria. Use suprema or infima unless attainment is established. An \(\varepsilon\)-optimal policy is feasible and within \(\varepsilon\) of the finite optimum value. Report an empty feasible set separately.

\textbf{Welfare and accounting.} Normative utility, interpersonal normalization, social weights, numeraire, discounting and the use of public receipts are primitives. Behavioral utility need not coincide with the welfare criterion. Household welfare includes ownership income and rents once; adding profits again to compensating variation generally double-counts them. A monetary equivalent is defined by an explicit transfer and price convention, with monotonicity and range conditions. Average effects, mechanism contributions, transfers and welfare losses are different targets. Infinite sums require convergence and dynamic feasible plans require terminal or no-Ponzi conditions.

\textbf{Source versions.} A working-paper date, revision date and journal publication year are distinguished. A DOI for a journal version is not automatically the DOI of an earlier manuscript. Locators refer to the stated version and printed pages where specified. A metadata-only check does not verify a claim about a full-text conclusion. Every entry's source subsection narrows the provenance claim accordingly.
% END ECONOMICS STATEMENT
\end{document}
